\documentclass[12pt,a4paper]{article}
\usepackage{amsmath}
\usepackage{amsfonts}
\usepackage{tgpagella}
\usepackage{graphicx}
\usepackage{polski}
\usepackage[utf8]{inputenc}
\usepackage{listings}
\usepackage{hyperref}

\pagestyle{empty}

\DeclareMathOperator{\sgn}{sgn}

\begin{document}

\title{Optymalizacja Kodu Na Różne Architektury\\Ćwiczenie 3}
\date{}

\maketitle


\section{Wprowadzenie}
Rozważmy następującą procedurę uruchamiającą eliminację Gaussa bez pivotingu

\lstinputlisting[language=C, frame=none, escapeinside={(*}{*)}, caption={Bazowa wersja kodu}, label={kod}, basicstyle=\small]{ge1.c}

\section{Ćwiczenia}

Przy kolejnych krokach proszę wykonywać kopie aktualnych wersji stanu kodu.
Będą potrzebne przy ostatnim punkcie oraz na kolejnych zajęciach.
Proszę na bieżąco notować czasy wykonywania się obliczeń.

\subsection{}


Proszę wprowadzić zaznaczone w kodzie modyfikacje a następnie skompilować procedurę bez optymalizacji.
Dla zadanego rozmiaru macierzy w kolejnych uruchomieniach wartość zmiennej \textbf{check} nie powinna ulegać zmianie.
\begin{lstlisting}[language=bash, frame=none, basicstyle=\small, escapeinside={(*}{*)}]
maciekw@Banach:~lab3> gcc ge1.c 
maciekw@Banach:~lab3> ./a.out 1500
call GETime: 1.314771e+00
FLOP: 3.371626e+09
GFLOP/s: 2.564421e+00
Check: -4.605781e+16
maciekw@Banach:~/studenci/OORA/skrypty/lab3>  
\end{lstlisting}

\subsection{}

Proszę wykonać kopię procedury
\begin{lstlisting}[language=bash, frame=none, basicstyle=\small, escapeinside={(*}{*)}]
maciekw@Banach:~lab3> cp ge1.c ge2.c
\end{lstlisting}
Proszę umieścić w rejestrach liczniki pętli.\\
Proszę skompilować procedurę z optymalizacją.\\
Jaki jest czas działania?\\
Jakie GFLOP/s uzyskuje procesor?

\subsection{}

Proszę wykonać kopię procedury
\begin{lstlisting}[language=bash, frame=none, basicstyle=\small, escapeinside={(*}{*)}]
maciekw@Banach:~lab3> cp ge2.c ge3.c
\end{lstlisting}
Proszę umieścić w rejestrze powtarzającą się w najbardziej zagnieżdżonym miejscu wartość.
\begin{lstlisting}[language=C, frame=none, basicstyle=\small, escapeinside={(*}{*)}]
multiplier = (A[i][k]/A[k][k]);
\end{lstlisting}
Proszę skompilować procedurę z optymalizacją.\\
Jaki jest czas działania?\\
Jakie GFLOP/s uzyskuje procesor?

\subsection{}

Proszę wykonać kopię procedury
\begin{lstlisting}[language=bash, frame=none, basicstyle=\small, escapeinside={(*}{*)}]
maciekw@Banach:~lab3> cp ge3.c ge4.c
\end{lstlisting}
Proszę ręcznie rozwinąć najbardziej zagnieżdżoną pętlę do 8 iteracji.
\begin{lstlisting}[language=C, frame=none, basicstyle=\small, escapeinside={(*}{*)}]
for (j = k+1; j < SIZE; ) {
  if (j < (max(SIZE - BLKSIZE, 0))) {
    //PROSZE UZUPELNIC
    j += BLKSIZE;
  }
  else{
    A[i][j]=A[i][j]-A[k][j]*multiplier;
    j++;
  }
\end{lstlisting}
Proszę skompilować procedurę z optymalizacją.\\
Jaki jest czas działania?

\subsection{}

Proszę wykonać kopię procedury
\begin{lstlisting}[language=bash, frame=none, basicstyle=\small, escapeinside={(*}{*)}]
maciekw@Banach:~lab3> cp ge4.c ge5.c
\end{lstlisting}
Proszę zamienić deklarację macierzy na macierz jednowymiarową indeksowaną poprzez makro (patrz laboratorium 2).\\
Proszę skompilować procedurę z optymalizacją.\\
Jaki jest czas działania?\\
Jakie GFLOP/s uzyskuje procesor?


%%%%%%%%%%%%%%%%%%%%%%%%%%%%%%5

\subsection{}

Proszę wykonać kopię procedury

\begin{lstlisting}[language=bash, frame=none, basicstyle=\small]
maciekw@Banach:~lab3> cp ge5.c ge6.c
\end{lstlisting}

Proszę wprowadzić blokowanie pamięci podręcznej oraz pakowanie danych do małych buforów roboczych.
Całość wykonujemy analogicznie do \textbf{\href{https://github.com/flame/how-to-optimize-gemm}{https://github.com/flame/how-to-optimize-gemm}}\\

\begin{enumerate}
  \item blokowanie zakresu aktualizowanej macierzy,
  \item pakowanie intensywnie używanych danych do buforów ciągłych,
  \item aktualizacja bloku za pomocą małego kernela obliczeniowego.
\end{enumerate}


Proszę dodać następujące definicje rozmiarów bloków:

\begin{lstlisting}[language=C, frame=none, basicstyle=\small]
#define MR 8
#define NR 8

#define MC 128
#define NC 256
\end{lstlisting}

Znaczenie parametrów jest następujące:

\begin{itemize}
  \item \textbf{NC} -- liczba kolumn w pakowanym fragmencie wiersza pivota,
  \item \textbf{MC} -- liczba wierszy aktualizowanych w jednym bloku,
  \item \textbf{NR} -- liczba kolumn przetwarzanych jednocześnie w ręcznie
        rozwiniętym kernelu.
\end{itemize}

Następnie proszę utworzyć funkcję pakującą fragment wiersza pivota:

\begin{lstlisting}[language=C, frame=none, basicstyle=\small]
static void PackPivotRow(int n, const double *src, double *dst)
{
  int j;

  for (j = 0; j < n; j++) {
    dst[j] = src[j];
  }
}
\end{lstlisting}

Funkcja ta kopiuje ciągły fragment wiersza \textbf{A[k][j]} do lokalnego bufora roboczego.
Spakowany fragment powinien być następnie użyty wiele razy dla różnych bloków wierszy.

Proszę utworzyć funkcję pakującą mnożniki eliminacji:

\begin{lstlisting}[language=C, frame=none, basicstyle=\small]
static void PackMultipliers(int m, const double *A, int SIZE,
                            int row0, int k, double pivot,
                            double *mult)
{
  int i;

  for (i = 0; i < m; i++) {
    mult[i] = A[IDX(row0 + i, k, SIZE)] / pivot;
  }
}
\end{lstlisting}

Funkcja ta oblicza wartości:

\begin{lstlisting}[language=C, frame=none, basicstyle=\small]
A[i][k] / A[k][k]
\end{lstlisting}

dla małego bloku wierszy i zapisuje je w lokalnym buforze.

Następnie proszę utworzyć kernel aktualizujący blok macierzy:

\begin{lstlisting}[language=C, frame=none, basicstyle=\small]
static void UpdateKernel(int m, int n, double *C,
                         int ldc, const double *pivot_row,
                         const double *mult)
{
  int i, j;

  for (j = 0; j + NR <= n; j += NR) {
    register double p0 = pivot_row[j];
    register double p1 = pivot_row[j + 1];
    register double p2 = pivot_row[j + 2];
    register double p3 = pivot_row[j + 3];
    register double p4 = pivot_row[j + 4];
    register double p5 = pivot_row[j + 5];
    register double p6 = pivot_row[j + 6];
    register double p7 = pivot_row[j + 7];

    for (i = 0; i < m; i++) {
      double *c = C + i * ldc + j;
      register double mi = mult[i];

      c[0] = c[0] - p0 * mi;
      c[1] = c[1] - p1 * mi;
      c[2] = c[2] - p2 * mi;
      c[3] = c[3] - p3 * mi;
      c[4] = c[4] - p4 * mi;
      c[5] = c[5] - p5 * mi;
      c[6] = c[6] - p6 * mi;
      c[7] = c[7] - p7 * mi;
    }
  }

  for (; j < n; j++) {
    register double pj = pivot_row[j];

    for (i = 0; i < m; i++) {
      C[i * ldc + j] = C[i * ldc + j] - pj * mult[i];
    }
  }
}
\end{lstlisting}

Po dodaniu funkcji pomocniczych należy zmienić funkcję \textbf{ge} na
następujący schemat:

\begin{lstlisting}[language=C, frame=none, basicstyle=\small]
int ge(double *A, int SIZE)
{
  int k;
  int ii, jj;

  double packed_pivot[NC];
  double packed_mult[MC];

  for (k = 0; k < SIZE; k++) {
    double pivot = A[IDX(k, k, SIZE)];

    for (jj = k + 1; jj < SIZE; jj += NC) {
      int jb = MIN(NC, SIZE - jj);

      PackPivotRow(jb, &A[IDX(k, jj, SIZE)], packed_pivot);

      for (ii = k + 1; ii < SIZE; ii += MC) {
        int ib = MIN(MC, SIZE - ii);

        PackMultipliers(ib, A, SIZE, ii, k, pivot, packed_mult);

        UpdateKernel(ib, jb, &A[IDX(ii, jj, SIZE)],
                     SIZE, packed_pivot, packed_mult);
      }
    }
  }

  return 0;
}
\end{lstlisting}



Proszę skompilować procedurę z optymalizacją:
Jaki jest czas działania?\\
Jakie GFLOP/s uzyskuje procesor?


%%%%%%%%%%%%



\subsection{}

Proszę wykonać kopię procedury
\begin{lstlisting}[language=bash, frame=none, basicstyle=\small, escapeinside={(*}{*)}]
maciekw@Banach:~lab3> cp ge6.c ge7.c
\end{lstlisting}
Proszę wprowadzić operacje wetorowe SSE3 (patrz laboratorium 2).
\begin{lstlisting}[language=C, frame=none, basicstyle=\small, escapeinside={(*}{*)}]
#include <x86intrin.h>
...
  for (j = 0; j + NR <= n; j += NR) {
    register __m128d p01 = _mm_loadu_pd(pivot_row + j);
    register __m128d p23 = _mm_loadu_pd(pivot_row + j + 2);
    register __m128d p45 = _mm_loadu_pd(pivot_row + j + 4);
    register __m128d p67 = _mm_loadu_pd(pivot_row + j + 6);

    for (i = 0; i < m; i++) {
      double *c = C + i * ldc + j;

      register __m128d mmult = _mm_set1_pd(mult[i]);

      register __m128d c01 = _mm_loadu_pd(c);
      register __m128d c23 = _mm_loadu_pd(c + 2);
      register __m128d c45 = _mm_loadu_pd(c + 4);
      register __m128d c67 = _mm_loadu_pd(c + 6);

      c01 = _mm_sub_pd(c01, _mm_mul_pd(p01, mmult));
      c23 = _mm_sub_pd(c23, _mm_mul_pd(p23, mmult));
      c45 = _mm_sub_pd(c45, _mm_mul_pd(p45, mmult));
      c67 = _mm_sub_pd(c67, _mm_mul_pd(p67, mmult));

      _mm_storeu_pd(c,     c01);
      _mm_storeu_pd(c + 2, c23);
      _mm_storeu_pd(c + 4, c45);
      _mm_storeu_pd(c + 6, c67);
    }
  }
\end{lstlisting}
Proszę skompilować procedurę z optymalizacją.\\
Jaki jest czas działania?\\
Jakie GFLOP/s uzyskuje procesor?

\subsection{}

Proszę wykonać kopię procedury
\begin{lstlisting}[language=bash, frame=none, basicstyle=\small, escapeinside={(*}{*)}]
maciekw@Banach:~lab3> cp ge7.c ge8.c
\end{lstlisting}
Proszę wprowadzić 256-bitowe operacje wetorowe AVX (patrz laboratorium 2).
\begin{lstlisting}[language=C, frame=none, basicstyle=\small, escapeinside={(*}{*)}]
#include <immintrin.h>
register __m256d ...
tmp0 = _mm256_loadu_pd(...));
...
\end{lstlisting}
Proszę skompilować procedurę z optymalizacją.
\begin{lstlisting}[language=bash, frame=none, basicstyle=\small, escapeinside={(*}{*)}]
maciekw@Banach:~lab3> gcc -mavx ge8.c
\end{lstlisting}
Jaki jest czas działania?\\
Jakie GFLOP/s uzyskuje procesor?


\subsection{}

Proszę wszystkie uzyskane na Laboratoriach kody skompilować z opcją optymalizacji \textbf{-O2} oraz przetestować opcje \textbf{-march=native} oraz \textbf{-mfma}.
\begin{lstlisting}[language=bash, frame=none, basicstyle=\footnotesize, escapeinside={(*}{*)}]
maciekw@Banach:~/studenci/OORA/skrypty/lab3> gcc -mavx -O2 ge7.c 
maciekw@Banach:~/studenci/OORA/skrypty/lab3> ./a.out 1500
call GETime: 3.592460e-01
FLOP: 3.371626e+09
GFLOP/s: 9.385284e+00
Check: -4.605781e+16
maciekw@Banach:~lab3> gcc -mavx -O2 -march=native ge8.c
maciekw@Banach:~lab3> ./a.out 1500
call GETime: 4.534000e-01
FLOP: 3.371626e+09
GFLOP/s: 7.436316e+00
Check: -4.605781e+16
maciekw@Banach:~lab3> gcc -mavx -O2 -mfma ge8.c
maciekw@Banach:~lab3> ./a.out 1500
call GETime: 2.618760e-01
FLOP: 3.371626e+09
GFLOP/s: 1.287489e+01
Check: -4.605781e+16
maciekw@Banach:~lab3> gcc -mavx -O2 -march=native -mfma ge8.c
maciekw@Banach:~lab3> ./a.out 1500
call GETime: 2.518740e-01
FLOP: 3.371626e+09
GFLOP/s: 1.338616e+01
Check: -4.605781e+16
maciekw@Banach:~lab3>  
\end{lstlisting}
Jaki jest ich czas działania?




\section{Podsumowanie}

Które metody optymalizacji dały najlepsze rezultaty?\\
Dlaczego? Jakie mechanizmy za nimi stały?


\end{document}
 

